% Alignment Forum post — ready to post 2026-04-29.
% Email drafts (TurboQuant authors) sent; Garg ping dropped.
% Compile: pdflatex day7_externalization.tex

\documentclass[11pt]{article}
\usepackage{amsmath,amssymb,amsthm}
\usepackage{microtype}
\usepackage{hyperref}
\usepackage[margin=1in]{geometry}
\usepackage{enumitem}
\usepackage{xcolor}
\usepackage{framed}

\newcommand{\R}{\mathbb{R}}
\newcommand{\E}{\mathbb{E}}
\newcommand{\cheb}{R_{\mathrm{c}}}

\title{Alignment Forum research note}
\author{Sankalp Pathak}
\date{2026-04-29}

\begin{document}
\maketitle

\begin{framed}\small
\noindent\textbf{Title:} Rate-distortion lower bound for LoRA model merging (research note, seeking feedback) \\
\textbf{Tags:} AI, Information Theory, Machine Learning \\
\textbf{Post type:} Research note / Draft
\end{framed}

\medskip

\noindent\textbf{TL;DR.} I'm a solo researcher working on a theoretical
paper about LoRA merging: \emph{how many bits do you need to store a
merged model that preserves $\epsilon$-distortion on each of $T$
tasks?} I have a 6-page note sketching a Shannon-style lower bound
that reduces to
\href{https://arxiv.org/abs/2504.19874}{TurboQuant Theorem~3} via a
deterministic $\max\geq\mathrm{avg}$ identity plus a rotationally-
invariant rate--distortion lemma on the task-vector centroid. Looking
for a sanity check on the reduction and on a Chebyshev-based
tail-event argument in Lemma~2.

\subsection*{Setting}

$T$ task vectors $\tau_1, \dots, \tau_T \in \R^d$ with $\|\tau_t\|
\leq B$. Per-task quadratic loss $f_t(w) = \tfrac12 \|w - \tau_t\|^2$.
A rate-$R$ merging algorithm is an encoder--decoder pair
$E: (\R^d)^T \to \{1, \dots, 2^R\}$, $D: \{1, \dots, 2^R\} \to \R^d$
producing $w^\star = D(E(\boldsymbol\tau))$. Figure of merit:
max-distortion $\Delta(w^\star; \boldsymbol\tau) = \max_t \|w^\star
- \tau_t\|^2$.

\subsection*{Result}

\[
  \mathcal{D}^\star(T, d, B, R)
  \;\geq\;
  B^2\!\left(1 - \tfrac{1}{T}\right)
  \;+\; c_{\mathrm{TQ}} \cdot \tfrac{B^2}{T} \cdot 2^{-2R/d}.
\]

Two terms with distinct meanings:
\begin{itemize}[leftmargin=1.4em,itemsep=0.15em]
  \item $B^2(1 - 1/T)$ is a \textbf{Chebyshev-radius floor} --- the
  irreducible multi-task incompatibility. Exactly $\E[(1/T)\sum_t
  \|\tau_t - \bar\tau\|^2]$ under iid uniform-sphere task vectors;
  numerically within 1--2\% of $\E[\cheb^2]$.
  \item $c_{\mathrm{TQ}} (B^2/T) \cdot 2^{-2R/d}$ is the classical
  Shannon decay governed by the rate--distortion function of the
  centroid $\bar\tau$.
\end{itemize}

\subsection*{Proof in one paragraph}

Yao's minimax with hard distribution $P = \mathrm{Unif}(B S^{d-1})^
{\otimes T}$. The deterministic identity
\[
  \tfrac{1}{T}\sum_t \|w^\star - \tau_t\|^2
  \;=\; \|w^\star - \bar\tau\|^2
  \;+\; \tfrac{1}{T}\sum_t \|\tau_t - \bar\tau\|^2
\]
holds because the cross-term $\sum_t(\tau_t - \bar\tau) = 0$
\textbf{identically}. Take max on the LHS and expectations:
$\E[\max_t \|w^\star - \tau_t\|^2] \geq \E\|w^\star - \bar\tau\|^2
+ B^2(1 - 1/T)$. Now apply a rate--distortion lower bound to the
first term: under $P$, $\bar\tau$ is rotationally invariant with
$\E\|\bar\tau\|^2 = B^2/T$, so the scale-invariant bound
$\E[\|\hat Y - Y\|^2 / \|Y\|^2] \geq c_{\mathrm{TQ}} \cdot 2^{-2R/d}$
from TurboQuant Thm~3 transfers (chain rule + DPI + Jensen).
Converting to absolute scale uses a high-probability norm-event
$G_d$ with $\Pr(G_d^c) = O(d^{-1})$, absorbing a $(1 + o_d(1))$
factor into the constant.

\subsection*{Novelty claim}

Standard multi-source rate--distortion frameworks (multiple
descriptions, Heegard--Berger, common reconstruction) all use
\textbf{per-source} or \textbf{average} distortion; none treat
$\max_t D_t$ as the figure of merit. The $\max$ is precisely what
isolates the multi-task Chebyshev floor, and the $\max\geq\mathrm{avg}$
reduction is what separates this floor from the compression decay.
I believe this is the first clean rate--distortion statement for
max-distortion multi-source with a single-vector-RD reduction.

\subsection*{Numerics}

$T \in \{2, 4, 8\}$, $d \in \{128, 512, 2048\}$, $b = R/d \in \{1,
2, 3, 4, 6, 8, 12\}$, 30 trials. Three checks:
\begin{itemize}[leftmargin=1.4em,itemsep=0.15em]
  \item \textbf{Chebyshev-floor formula $B^2(1 - 1/T)$.}
  Empirical/theory ratio in $[0.9984, 1.0173]$ across all $(T, d)$.
  \item \textbf{LB respected.} Valid everywhere (as it must be).
  \item \textbf{$2^{-R/d}$ regime --- resolved NO.} A quantized-
  Chebyshev-center merge closes the 2-regime apparent decay; the
  $2^{-R/d}$ behavior in an earlier mean-merge experiment was a
  sub-optimal-centering plateau, not an intrinsic RD phenomenon. The
  rate--distortion function is $\Theta\bigl(\cheb^2 + (B^2/T)
  \cdot 2^{-2R/d}\bigr)$.
\end{itemize}

\subsection*{What I'm asking}

Two specific checks:
\begin{enumerate}[leftmargin=1.6em,itemsep=0.15em]
  \item \textbf{Lemma~2's tail-event substitution.} $Y = \bar\tau$
  has no deterministic lower bound on $\|Y\|^2$ ($T=2$, $\tau_1 =
  -\tau_2$ kills it). I use $G_d = \{\|\bar\tau\|^2 \geq
  B^2/(T\kappa_d)\}$ with $\Pr(G_d^c)$ controlled by
  $\mathrm{Var}(\|\bar\tau\|^2) = O(B^4/(T^2 d))$. Is there a reason
  to prefer an EPI-based derivation that bypasses the
  $\rho$-conditioning entirely?
  \item \textbf{Novelty.} I checked El Gamal--Cover 1982, Steinberg
  2009, Heegard--Berger, Wyner--Ziv; none match. Am I missing an
  earlier treatment of max-over-$T$ distortion that I should be
  citing?
\end{enumerate}

Any other pointers --- ``you missed reference $X$'', ``Lemma~2 is
cleaner if you do $Y$'', ``this is subsumed by $Z$'' --- very
welcome. Full 6-page note in comments / first reply (or DM me if
preferred).

\end{document}
